% !TEX root = ../../main.tex
\chapter{Introduction}

\section{Background}

The integration of \textit{inverter-based resources} (IBRs) into transmission systems changes the fundamental assumptions traditionally used in transmission line protection design. Unlike synchronous generators, which generally supply large and relatively predictable fault currents, IBR fault currents are limited by converters, control strategies, and interconnection requirements. Consequently, their magnitudes, phase angles, and sequence-component content during faults differ from those of conventional systems \cite{haddadi2021systemprotection,kou2020fault,ieee2022std2800}. This change is important because most transmission line protection schemes are still designed based on the assumption of a system dominated by synchronous machines.

In transmission line protection, distance relays (relay 21) and line current differential relays (relay 87L) are widely used protection functions because each offers good selectivity, speed, and coverage for transmission line faults \cite{ieee2016c37113,ieee2015c37243}. However, when IBRs are connected to the grid, the quantities on which these relays base their decisions do not always follow conventional system behavior. Relay 21 relies on the apparent impedance estimated from the voltage-to-current ratio. Therefore, changes in the IBR fault-current contribution may affect impedance measurement accuracy and relay operating decisions \cite{ieee2016c37113,hooshyar2015cirepp1,fang2019impact}. Meanwhile, relay 87L compares the currents at both ends of a transmission line. Changes in the amplitude, phase angle, and initial dynamics of fault currents caused by IBRs may therefore affect the relationship between the differential current and the \textit{restraint current} that defines its operating characteristic \cite{ieee2015c37243,li2017analysis,chowdhury2023lcd}.

These issues were comprehensively discussed by Chowdhury and Fischer in \textit{Transmission Line Protection for Systems With Inverter-Based Resources---Part I: Problems} \cite{chowdhury2021tlp1}. The authors demonstrated that IBR responses during faults challenge the reliability of conventional transmission line protection schemes. Their collaborative study used \textit{black-box} electromagnetic transient (EMT) models from several IBR manufacturers and included a mix of Type 3 wind, Type 4 wind, and solar photovoltaic (PV) resources \cite{chowdhury2021tlp1}. The results showed that replacing synchronous sources with IBRs can create protection operating conditions that no longer conform to conventional system expectations. In \textit{Part II}, Chowdhury and Fischer demonstrated that specific setting modifications could substantially improve reliability \cite{chowdhury2021tlp2}. This finding confirms that the primary issue is not merely the presence of IBRs, but the mismatch between new system behavior and legacy protection settings.

With specific regard to relay 21, Hooshyar \textit{et al.} showed that lines connected to fully converter-interfaced renewable power plants may experience failure to detect internal faults or undesired tripping for external faults \cite{hooshyar2015cirepp1}. Fang \textit{et al.} also demonstrated that \textit{inverter-interfaced renewable energy generators} can affect distance relay performance, indicating the need for enhanced protection schemes \cite{fang2019impact}. Regarding relay 87L, Li \textit{et al.} found that inverter-based renewable generation can alter the relationship between operating and restraint currents in line current differential protection \cite{li2017analysis}. Chowdhury \textit{et al.} subsequently showed that, in systems with IBRs, the dependability of 87L elements may decrease and their security may deteriorate if they are not properly applied \cite{chowdhury2023lcd}. Therefore, both relay 21 and relay 87L warrant further investigation in transmission systems with IBR penetration.

Based on these conditions, this research is designed to replicate the problem framework presented in \cite{chowdhury2021tlp1}, while narrowing the focus to two protection functions: relay 21 and relay 87L. Rather than examining all transmission line protection schemes considered in the reference study, this research specifically investigates how the characteristics of these two relays change when IBRs are introduced into a transmission system that was initially assumed to be conventional. The primary variable is the IBR type, namely Type 3 wind turbines, Type 4 wind turbines, and photovoltaic systems, while the protection settings remain based on the conventional system. Through this approach, the research is expected to identify the applicability limits of conventional settings and provide a basis for developing technical solutions for transmission line protection in systems with IBRs.

\section{Problem Formulation}

Based on the background described above, the research problems are formulated as follows:

\begin{enumerate}
	\item The operating characteristics of relay 21 in a transmission system integrated with \textit{inverter-based resources} (IBRs) when conventional protection settings are retained.
	\item The operating characteristics of relay 87L in a transmission system integrated with IBRs when conventional protection settings are retained.
	\item The effects of different IBR types, namely Type 3 wind turbines, Type 4 wind turbines, and photovoltaic systems, on the responses of relay 21 and relay 87L.
	\item The suitability of conventional protection settings for transmission systems integrated with IBRs.

\end{enumerate}

\section{Research Objectives}

The objectives of this research are as follows:

\begin{enumerate}
	\item To analyze the operating characteristics of relay 21 in a transmission system with IBR integration using conventional protection settings.
	\item To analyze the operating characteristics of relay 87L in a transmission system with IBR integration using conventional protection settings.
	\item To compare the effects of Type 3 wind turbines, Type 4 wind turbines, and photovoltaic systems on the responses of relay 21 and relay 87L.
	\item To evaluate the suitability of conventional protection settings for transmission systems containing IBRs.
	      % \item To formulate proposed technical solutions concerning the applicability limits of conventional settings,
	      %       required setting adjustments, and potential supplementary protection approaches
	      %       based on the analysis results.
\end{enumerate}

\section{Research Scope and Limitations}

To maintain a clear research focus, the scope and limitations of this research are defined as follows:

\begin{enumerate}
	\item The research focuses on transmission systems with IBR integration.
	\item The analysis is limited to relay 21 and relay 87L.
	\item This research replicates the problem perspective presented in \cite{chowdhury2021tlp1}, but limits the investigation to relay 21 and relay 87L.
	\item The IBR types analyzed are Type 3 wind turbines, Type 4 wind turbines, and photovoltaic systems.
	\item The IBR penetration level is kept constant; therefore, the primary research variable is the IBR type.
	\item Each relay is evaluated individually; protection coordination and primary-backup protection schemes are not considered.
	\item The fault scenarios include faults on the IBR-connected line, faults on another line not connected to the IBR, and a condition in which one transmission line is \textit{out of service}.
	\item The research is conducted through simulations using DIgSILENT PowerFactory.
\end{enumerate}

\section{Research Contributions}

This research is expected to provide the following contributions:

\begin{enumerate}
	\item \textbf{Academic contributions}
	      \begin{itemize}
		      \item To expand the literature on the effects of IBRs on transmission line protection, particularly relay 21 and relay 87L.
		      \item To provide a basis for further research on protection setting adjustments or adaptive protection in transmission systems with IBRs.
	      \end{itemize}
	\item \textbf{Practical contributions}
	      \begin{itemize}
		      \item To provide insight into the conditions under which conventional protection settings remain suitable for systems with IBRs.
		      \item To provide technical input regarding the need to reassess relay 21 and relay 87L settings in transmission systems with IBRs.
		      \item To provide a basis for formulating protection recommendations for transmission systems beginning to integrate Type 3 wind turbines, Type 4 wind turbines, and photovoltaic systems.
	      \end{itemize}
\end{enumerate}

\section{Thesis Organization}

This thesis is organized as follows.

\noindent Chapter I presents the introduction, which consists of the background, problem formulation, research objectives, research scope and limitations, research contributions, and thesis organization.

\noindent Chapter II presents the literature review and theoretical foundations, covering conventional power systems and IBR integration; the fault characteristics of Type 3 wind turbines, Type 4 wind turbines, and photovoltaic systems; the operating principles of relay 21 and relay 87L; the effects of IBRs on transmission protection performance; a review of previous studies; and the position and contribution of this research.

\noindent Chapter III presents the research methodology, including the system model, study assumptions, modeling of conventional sources and IBRs, modeling of relay 21 and relay 87L, test scenarios, evaluation parameters, simulation procedures, and methods used to analyze the results.

\noindent Chapter IV presents the research results and discussion, including the baseline system condition and initial verification; a comparison of short-circuit current contributions among IBR types; the decomposition of current sequence components for each scenario; the responses of relay 21 and relay 87L in systems with Type 3 wind turbines, Type 4 wind turbines, and photovoltaic systems; a comparison of the characteristics of both relays among the IBR types; the implications of retaining conventional settings; and a summary of relay suitability in IBR-integrated systems.

\noindent Chapter V presents the conclusions and recommendations. The conclusions answer the research problems and objectives based on the findings, while the recommendations identify directions for further research that are not covered within the scope and limitations of this study.
